\label{sec:related_work}

\textbf{Ontology-based IR and food ontologies.}
Classical ontology-based retrieval enriches document and query representations with explicit semantic structure to reduce vocabulary mismatch and improve precision \cite{vallet2005ontology, castells2007adaptation, fernandez2011semantic}. Subsequent studies have further explored the integration of semantic knowledge graphs and structural constraints into information retrieval and cross-modal search architectures \cite{chen2021personalized, carvalho2018crossmodal, lien2020recipe}. In the food domain, several comprehensive food ontologies and knowledge bases have been constructed to formalize culinary concepts, traceability, and ingredient relationships \cite{dooley2018foodon, haussmann2019foodkg, recipemind2022}. However, these structured food ontologies are predominantly leveraged for downstream tasks such as personalized recommendation or ingredient substitution \cite{fatemi2023learning}, rather than direct ranking-based information retrieval. Crucially, neither line of research integrates ontology with modern neural retrieval pipelines for low-resource languages. To the best of our knowledge, there is currently no existing literature that tackles Vietnamese food information retrieval using an ontology-based approach. Consequently, our work stands as one of the pioneering efforts in the field of Vietnamese food IR, explicitly emphasizing the use of a localized food ontology to enhance dense retrieval pipelines and resolve unique culinary semantics in the Vietnamese language.

\textbf{Food retrieval under neural frameworks.}
Recent work increasingly treats food search as a retrieval problem within neural pipelines. Early recipe retrieval studies focused on cross-modal alignment between recipes and images \cite{carvalho2018crossmodal, lien2020recipe}. More recent benchmarks such as Recipe-MPR \cite{zhang2023recipempr} highlight the difficulty of multi-attribute preference-based retrieval, while Hu et al.~\cite{hu2024bridging} extend recipe retrieval to cross-cultural settings through reformulation and reranking. KERL \cite{mohbat2025kerl} further shows that structured food knowledge remains useful even in LLM-augmented pipelines. In the broader IR community, dense retrieval has become the dominant paradigm for semantic matching \cite{karpukhin2020dense, thakur2021beir}, and recent graph-aware retrieval frameworks demonstrate that explicit relational structure can complement purely embedding-based retrieval on compositional queries \cite{he2024g, hipporag2024, kg2rag2025}. Aligning with this prominent trend, our proposed framework harnesses the power of neural pipelines by adopting dense retrieval as the core semantic matching engine. By injecting an explicitly constructed Vietnamese food ontology into this neural architecture, we bridge the gap between advanced representation learning and explicit relational reasoning, effectively combining robust semantic embeddings with the structural knowledge necessary for food-specific retrieval tasks.

\textbf{Dish similarity and related-dish recommendation.}
While some prior work has addressed ingredient-level substitution using knowledge graphs \cite{shirai2021identifying} or learning-based models \cite{fatemi2023learning}, our focus shifts toward the broader challenge of modeling holistic dish similarity. Existing structured approaches often capture relations through food-chemical representations such as FlavorGraph \cite{park2021flavorgraph}. However, lightweight baselines for related-dish recommendation typically compute similarity based on flat ingredient overlap. This flat approach is fundamentally limited as it cannot distinguish between substitutions within the same semantic class (e.g., swapping two types of meat) and cross-class differences. To directly support our related-dish recommendation task, we address this gap by moving beyond simple ingredient overlap, leveraging our ontology to compute a hierarchy-aware similarity that incorporates structural category overlap, cooking methods, and flavor profiles.

\textbf{Position of the present work.}
Our work bridges these three lines by: (1)~constructing a Vietnamese-specific ontology from a 10K-dish dataset, (2)~integrating it into a dense-retrieval pipeline as a pre-embedding enrichment mechanism for two food retrieval tasks, and (3)~providing an explicit ablation that isolates the contribution of hierarchy, typed relations, and inference rules.
